\begin{abstract}
Бинарные структуры из коллоидной квантовой точки (КТ) и металлической
наночастицы (МНЧ) позволяют управлять возбуждением экситона за счёт
плазмонного усиления локального поля. Исследовано резонансное возбуждение
двухуровневой КТ гауссовым лазерным импульсом длительностью
\SI{20}{\femto\second} вблизи золотого вытянутого сфероида, приближающего
форму наноцилиндра. Использована полуклассическая модель, объединяющая
уравнения Блоха и причинный дисперсионный отклик МНЧ. Раздельно сопоставлены
диполь-дипольное взаимодействие и полное квазистатическое поле сфероида,
а также одноосцилляторное и многоосцилляторное описания поляризуемости
золота. Рассмотрены размещение КТ у вершины и сбоку, продольная и поперечная
поляризации возбуждающего поля. Для выбранных параметров наименьший флюенс,
необходимый для достижения населённости возбуждённого уровня $0.5$,
получен при расположении КТ у вершины и поле вдоль оси МНЧ. При поверхностном
зазоре \SI{1}{\nano\meter} расчёт даёт снижение порога примерно в $38$ раз
относительно изолированной КТ; дипольная модель завышает этот порог в $2.4$
раза. В том же канале при зазоре \SI{0.5}{\nano\meter}
одноосцилляторная аппроксимация завышает порог в $7.7$ раза.
Оптимизация зазора зависит от схемы возбуждения:
слабополевой отклик на фиксированной частоте имеет внутренний максимум,
тогда как импульсный порог уменьшается до нижней границы исследованного
диапазона. Предложена схема различения преимущественных конфигураций
ансамбля по поляризационным спектрам, зависимостям возбуждения от флюенса
и кинетике люминесценции; рассмотрены доступные способы химической сборки
и ориентации частиц. Количественные рекомендации условны относительно
точечного описания КТ, квазистатики и принятых скоростей релаксации;
изменение спонтанного распада и квантовый выход люминесценции отдельно
не рассчитываются.
\end{abstract}

\noindent\textbf{Ключевые слова:} коллоидная квантовая точка, золотой
наностержень, сфероид, экситон-плазмонное взаимодействие, фемтосекундный
импульс, поляризационная спектроскопия.
